\xiaosi

NCP的布线拓扑由6个离散参数决定：中间神经元数$n_i$、命令神经元数$n_c$、感觉扇出$f_s$、中间扇出$f_i$、命令层循环连接数$r_c$和运动扇入$f_m$，参数之间还存在不等式约束（例如$f_s \leq n_i$）。这类离散、带约束的组合优化问题难以用基于梯度的方法处理，因此本文选择进化搜索算法\cite{stanley2002neat}来探索布线空间。

进化搜索模拟自然选择的过程，基本流程为：随机初始化一个由若干布线配置组成的种群，评估每个个体的适应度，然后通过选择、交叉、变异产生下一代，反复迭代直至收敛。本文采用的具体算子如下：

选择阶段使用锦标赛选择，即从种群中随机抽取$k$个个体，取适应度最高者作为父代。$k$值越大，选择压力越强，收敛越快但多样性损失也越大。交叉阶段对两个父代的6个基因逐位独立地以概率0.5选择来源，拼接为子代基因组。变异阶段对每个基因以概率$p_{mut}$施加高斯扰动$\Delta \sim \mathcal{N}(0, \sigma^2)$，取整后裁剪到该基因的合法范围内。此外，每代保留适应度最高的若干个体直接进入下一代（精英保留），防止已发现的优质解在随机操作中丢失。

需要注意的是，交叉和变异可能产生违反约束的基因组合。例如，若变异后$f_s$超过了$n_i$，则实际构建NCP布线时会出错。为此，本文在每次遗传操作后增加一步约束修复，将违约参数裁剪到合法范围。这一设计参考了带约束进化优化的常见做法。

相比可微分架构搜索方法（如DARTS\cite{liu2019darts}需要将离散选择松弛为连续权重），进化方法直接在离散空间上操作，不需要对搜索空间做连续化近似，实现也更为简单。
